% `tex.web` §1223: `\meaning` distinguishes a `\chardef` from a `\mathchardef`.
% They are printed from different `chr_cmd` strings — `char_given` is `\char`
% and `math_given` is `\mathchar` — because they are different things: a
% `\chardef`'d constant is a character to typeset, a `\mathchardef`'d one is an
% atom carrying a class, a family and a position.
%
% We printed `\char` for both. Every symbol plain.tex defines — `\alpha` is
% `\mathchardef"010B` — then claimed through `\meaning` and `\show` to be an
% ordinary character constant, and a macro switching on `\meaning` could not
% tell a relation from a letter.
%
% Found in trip.tex, where `\show` of a `\mathchardef`'d `\minus` reads
% `\mathchar"232D` in Knuth's transcript and read `\char"232D` in ours.
\catcode`\{=1 \catcode`\}=2

\mathchardef\mc="232D
\chardef\ch="41
\message{<mc=\meaning\mc>}
\message{<ch=\meaning\ch>}

% `\the` is the same integer for both — only the *meaning* distinguishes them,
% which is the whole point of the two strings.
\message{<mcval=\the\mc><chval=\the\ch>}

% a mathchar big enough to need all four hex digits, and one small enough not to
\mathchardef\big="7FFF
\mathchardef\small="1
\message{<big=\meaning\big><small=\meaning\small>}

% and `\ifx` sees the two as different meanings even at the same value — which
% has to be tested inside 0..255, since a character code wider than that is
% §434's error and not a `\chardef` at all.
\mathchardef\msame="41
\message{<ifx=\ifx\ch\msame yes\else no\fi>}
\message{<ifxself=\ifx\ch\ch yes\else no\fi>}
\end
